\section*{Capítulo \ref{ch:b1:radiation-heat-transport} --- Radiação e transporte de calor}
% values: calc/ch27.py

\begin{solution}{exo:b1:radiation-heat-transport:1}
(a)~$S_0/\pi = 1361/\pi = \qty{433.2}{W.m^{-2}}$. (b)~No equador num equinócio $\bar Q = (S_0/\pi)(a/r)^2$,
logo $(a/r)^2 = 436.7/433.2 = 1.0080$ e $r/a =
1.0080^{-1/2} = 0.9960$: a Terra está \qty{0.4}{\percent} mais perto do Sol que sua distância média.
(c)~$\bar Q = 433.2 \times 0.9683 \times \cos\ang{23.44} =
\qty{384.9}{W.m^{-2}}$.
\end{solution}

\begin{solution}{exo:b1:radiation-heat-transport:2}
(a)~$415.6/172.4 = 2.41$. (b)~A zona entre os paralelos $\pm\varphi$ tem a área $2\pi a_\oplus^2 \times 2\sin\varphi$, uma fração $\sin\varphi$ da esfera: $\sin\ang{37.0} = 0.602$, \qty{60}{\percent}.
(c)~Num eixo em \omterm{def:b1:spherical-astronomy:geographic}{latitude}, as faixas perto dos polos parecem tão largas quanto as perto do equador,
mas cobrem bem menos área: a zona a polo de \ang{60} é apenas $1 - \sin\ang{60} =
\qty{13.4}{\percent}$ da superfície. Contra $\sin\varphi$, os trópicos ocupam a maior parte do eixo.
\end{solution}

\begin{solution}{exo:b1:radiation-heat-transport:3}
(a)~$T_e = (\text{OLR}/\sigma)^{1/4}$: \qty{256.4}{K}, \qty{241.3}{K} e \qty{254.5}{K}.
(b)~Luz solar absorvida: $287.7/75.8 = 3.80$; OLR: $244.9/192.3 =
1.27$. (c)~A \omterm{def:b1:radiation-heat-transport:net-radiation}{radiação líquida} segue a luz solar absorvida: positiva e grande nos trópicos,
fortemente negativa nas altas \omterm{def:b1:spherical-astronomy:geographic}{latitudes}. A OLR, que varia pouco, não compensa a queda da luz solar
absorvida.
\end{solution}

\begin{solution}{exo:b1:radiation-heat-transport:4}
(a)~$2.82 \times 365.25 = \qty{1030}{mm}$, \qty{1.03}{m} por ano. (b)~$1.030 \times
78.1 = \qty{80.4}{W.m^{-2}}$.
(c)~O $LE$ da reanálise é \qty{2}{W.m^{-2}} maior. Ao longo de muitos anos a atmosfera não pode
continuar acumulando vapor d'água: toda a água evaporada volta a cair, de modo que $LE$ e $LP$ devem
ser iguais na média global. A diferença é um desequilíbrio do modelo, como seu desequilíbrio
radiativo.
\end{solution}

\begin{solution}{exo:b1:radiation-heat-transport:5}
(a)~$\bar Q = S_0(a/r)^2\sin\delta = 1361 \times 0.9683 \times 0.3978 =
\qty{524.2}{W.m^{-2}}$. (b)~$1361 \times 1.0334 \times 0.3978 =
\qty{559.5}{W.m^{-2}}$; a razão é $1.0334/0.9683 = 1.067$. (c)~O polo tem luz solar 24 horas em 24
a uma altitude constante, de modo que seu $\sin h_\odot$ médio no dia é $\sin\delta = 0.398$. No
equador o Sol fica alto ao meio-dia, mas abaixo do horizonte metade do tempo; a média de
$\sin h_\odot$ sobre as 24 horas é $\cos\delta/\pi =
0.292$. A razão é $\pi\tan\delta = 1.362$ (\cref{prop:b1:radiation-heat-transport:polar-summer}).
\end{solution}

\begin{solution}{exo:b1:radiation-heat-transport:6}
(a)~$\cos h_0 = -\tan\ang{40}\tan\ang{23.44} = -0.3638$, $h_0 = \ang{111.33} =
\qty{1.943}{rad}$; o dia dura $24 \times 111.33/180 = \qty{14.84}{h}$; $\bar Q =
(1361/\pi) \times 0.9683 \times (1.943 \times 0.6428 \times 0.3978 + 0.7660 \times
0.9175 \times 0.9315) = \qty{483.0}{W.m^{-2}}$. (b)~$\cos h_0 = +0.3638$, $h_0 =
\ang{68.67}$, \qty{9.16}{h}; $\bar Q = \qty{155.9}{W.m^{-2}}$. (c)~A razão é
$483.0/155.9 = 3.10$. A distância contribui com $0.9683/1.0334 = 0.937$: ela trabalha contra o
verão do norte. A geometria, a altitude do Sol e a duração do dia contribuem com $3.10/0.937 = 3.31$,
toda a diferença sazonal e mais.
\end{solution}

\begin{solution}{exo:b1:radiation-heat-transport:7}
(a)~$B = H/LE$: \num{0.093} no equador, \num{0.374} a \ang{25} N, \num{0.338} a \ang{45} N,
\num{-0.152} a \ang{70} N. (b)~\ang{25} N: a média zonal inclui os grandes desertos quentes, terra
seca com pouca água para evaporar, onde o solo aquecido aquece o ar em vez disso. (c)~$H < 0$
significa que o calor flui do ar para a superfície: na média anual a neve e o gelo marinho são mais
frios que o ar acima deles, que é aquecido por ar trazido de \omterm{def:b1:spherical-astronomy:geographic}{latitudes} mais baixas.
\end{solution}

\begin{solution}{exo:b1:radiation-heat-transport:8}
(a)~Avaliado: $\alpha = 100/340 = 0.294$, $R = 340 - 100 - 239 =
\qty{+1}{W.m^{-2}}$ (o \num{0.7} arredondado). Reanálise: $\alpha = 117.0/341.9 =
0.342$, $R = 341.9 - 117.0 - 237.7 = \qty{-12.8}{W.m^{-2}}$ com estes termos arredondados,
\qty{-12.9}{W.m^{-2}} com os não arredondados de
\cref{ex:b1:radiation-heat-transport:global}. (b)~A reanálise absorve $224.8$ contra $240$,
\qty{15.2}{W.m^{-2}} a menos, porque reflete \qty{17}{W.m^{-2}} a mais; sua OLR é apenas
\qty{1.3}{W.m^{-2}} menor. A reflexão em excesso explica o desequilíbrio. (c)~Integrando $R$ do Polo
Sul ao Polo Norte obtém-se a integral global, $-12.9 \times 4\pi a_\oplus^2 = -12.9 \times
5.10\times10^{14} = \qty{-6.6}{PW}$. Um transporte através do Polo Norte é impossível: não há nada
além dele, o paralelo é um ponto. A proposição vale apenas se $\int R\,\mathrm dA = 0$, o que o
regime permanente exige.
\end{solution}

\begin{solution}{exo:b1:radiation-heat-transport:9}
(a)~A área é $2\pi a_\oplus^2 \times 2\sin\ang{30} = 2\pi a_\oplus^2 =
\qty{2.55e14}{m^2}$, metade da Terra; o excedente é $33 \times 2.55\times10^{14} =
\qty{8.4}{PW}$. (b)~\qty{4.2}{PW} através de cada paralelo. (c)~Os dois somam \qty{8.5}{PW}, o
excedente, mas são desiguais: o hemisfério sul exporta \qty{0.5}{PW} a mais que o do norte. O
excedente não é dividido no equador: nesta reanálise \qty{0.36}{PW} do excedente do trópico norte
cruza o equador para o sul.
\end{solution}

\begin{solution}{exo:b1:radiation-heat-transport:10}
(a)~Escreva a média diária como $\bar Q = (a/r)^2 F(\varphi, \delta)$, com $F$ o colchete de
\cref{thm:b1:radiation-heat-transport:daily} vezes $S_0/\pi$. Durante $\mathrm dt$ a energia
recebida por metro quadrado é $\bar Q\,\mathrm dt$, e $\mathrm dt = r^2\,\mathrm d\lambda/h$, logo, então
$\bar Q\,\mathrm dt = (a^2/h)
F(\varphi, \delta(\lambda))\,\mathrm d\lambda$: o fator $r^2$ se cancela. (b)~Ao longo de um ano a
energia é $(a^2/h)\int_0^{2\pi} F(\varphi,
\delta(\lambda))\,\mathrm d\lambda$. A fórmula não muda quando $\varphi$ e $\delta$ mudam de sinal
juntos, $F(-\varphi, \delta) = F(\varphi, -\delta)$, e $\delta(\lambda + \pi) = -\delta(\lambda)$, já
que $\sin\delta = \sin\varepsilon
\sin\lambda$. Deslocando $\lambda$ de $\pi$ na integral, mostra-se que $-\varphi$ recebe a mesma
energia que $\varphi$. (c)~A primavera e o verão do norte duram \num{7.6} dias a mais, porque a Terra
está perto do afélio e se move devagar então; o verão austral é mais brilhante, mas mais curto. Pela
(a), os dois efeitos se cancelam exatamente.
\end{solution}

\begin{solution}{exo:b1:radiation-heat-transport:11}
(a)~$\langle (a/r)^2\rangle = \frac1P\int_0^P \frac{a^2}{r^2}\,\mathrm dt =
\frac1P\int_0^{2\pi}\frac{a^2}{r^2}\frac{r^2}{h}\,\mathrm d\lambda = \frac{2\pi
a^2}{hP} = \frac{2\pi a^2}{2\pi a b} = \frac ab$, e $b = a\sqrt{1-e^2}$.
(b)~A média global instantânea $\frac14 S_0(a/r)^2$ tem média $\frac14 S_0
\langle(a/r)^2\rangle = S_0/(4\sqrt{1-e^2})$. (c)~$340.25/\sqrt{1 - 0.0167^2} =
\qty{340.30}{W.m^{-2}}$, \qty{0.05}{W.m^{-2}} ou \qty{0.014}{\percent} acima de $S_0/4$.
\end{solution}

\begin{solution}{exo:b1:radiation-heat-transport:12}
(a)~Atmosfera $0.78 \times 4.10 = \qty{3.20}{PW}$, oceano $0.22 \times 4.10 =
\qty{0.90}{PW}$. (b)~A atmosfera sozinha transporta mais, \qty{5.0}{PW}, que a atmosfera e o oceano
juntos desta reanálise, \qty{4.02}{PW}: o transporte total da reanálise é menor em pelo menos um
quinto ($1 - 4.02/5.0 = 0.20$). Sua \omterm{def:b1:radiation-heat-transport:net-radiation}{radiação líquida} varia pouco demais entre os trópicos e os polos.
(c)~Remover uma constante de $R$ torna nula sua integral global, mas deixa inalterada sua variação
com a \omterm{def:b1:spherical-astronomy:geographic}{latitude}; o erro está nessa variação, no albedo das nuvens e superfícies do modelo em cada
\omterm{def:b1:spherical-astronomy:geographic}{latitude}, e não apenas na média global.
\end{solution}

\begin{solution}{pb:b1:radiation-heat-transport:1}
\begin{enumerate}
  \item Junho: $\cos h_0 = -0.3638$, $h_0 = \ang{111.33}$, $\bar Q =
        \qty{483.0}{W.m^{-2}}$. Dezembro: $h_0 = \ang{68.67}$, $\bar Q =
        \qty{155.9}{W.m^{-2}}$. Março: $h_0 = \ang{90}$, $\bar Q = (1361/\pi)
        \times 1.0081 \times \cos\ang{40} = \qty{334.6}{W.m^{-2}}$.
  \item Junho: $\tan\ang{80}\tan\ang{23.44} = 2.46 > 1$, dia polar, $h_0 = \pi$:
        $\bar Q = 1361 \times 0.9683 \times \sin\ang{80}\sin\ang{23.44} =
        \qty{516.3}{W.m^{-2}}$. Dezembro: noite polar, $\bar Q = 0$. Março:
        $(1361/\pi) \times 1.0081 \times \cos\ang{80} = \qty{75.8}{W.m^{-2}}$.
  \item A \ang{80} N o Sol não se põe ($h_0 = \ang{180}$, 24 horas), enquanto a \ang{40} N ele
        brilha por \qty{14.8}{h}. Sua altitude é menor, entre \ang{13.4} e \ang{33.4}, mas o termo
        $h_0\sin\varphi\sin\delta$, que cresce com a duração do dia e a \omterm{def:b1:spherical-astronomy:geographic}{latitude}, prevalece.
  \item As médias das faixas excedem os valores centrais em \qty{0.5}{\percent},
        \qty{0.8}{\percent}, \qty{1.7}{\percent}, \qty{1.7}{\percent} e \qty{1.3}{\percent}.
        Primeiro, a reanálise usa uma irradiância solar ligeiramente maior (Questão~5),
        \qty{0.5}{\percent} a mais. Segundo, a média de uma faixa não é o valor no centro: a metade
        equatorial de cada faixa tem mais área e mais luz solar que a metade polar.
  \item Pela \cref{prop:b1:radiation-heat-transport:global}, $S_0 = 4\sqrt{1-e^2}
        \times 341.9 = \qty{1367}{W.m^{-2}}$, \qty{0.47}{\percent} a mais que \qty{1361}{W.m^{-2}}.
  \item $A = 2\pi a_\oplus^2(1 - \sin\ang{40}) = 2\pi \times (6.371\times10^6)^2
        \times 0.3572 = \qty{9.110e13}{m^2}$, uma fração $(1 - \sin\ang{40})/2 =
        0.179$ da Terra.
  \item $(1 - \alpha)\bar Q$: \num{198.0}, \num{159.3}, \num{124.7}, \num{92.5},
        \num{74.7} \unit{W.m^{-2}}.
  \item $R$: \num{-28.1}, \num{-51.0}, \num{-76.6}, \num{-104.9}, \num{-116.9}
        \unit{W.m^{-2}}. Todas as cinco faixas estão em déficit; a faixa polar \ang{80}--\ang{90} tem o maior.
  \item $T_e = (\text{OLR}/\sigma)^{1/4}$: \num{251.3}, \num{246.8}, \num{244.1},
        \num{242.9}, \num{241.1} \unit{K}.
  \item Absorvida: $198.0/74.7 = 2.65$; OLR: $226.1/191.6 = 1.18$.
  \item O Sol está mais baixo, e a luz que chega com ângulo rasante é mais refletida pelas nuvens e
        pela água; neve e gelo marinho, superfícies brilhantes, cobrem mais do solo e do mar em
        direção ao polo.
  \item A OLR varia como $T_e^4$, e $T_e$ varia apenas \qty{10}{K}, \qty{4}{\percent}, pela calota:
        as altas \omterm{def:b1:spherical-astronomy:geographic}{latitudes} são mantidas bem mais quentes do que sua luz solar sozinha permitiria,
        pelo calor trazido do sul.
  \item Em \unit{m^2}: \num{3.143e13}, \num{2.550e13}, \num{1.879e13}, \num{1.151e13},
        \num{0.387e13}; a soma é \qty{9.110e13}{m^2}.
  \item $R_iA_i$: \num{-0.883}, \num{-1.300}, \num{-1.439}, \num{-1.207}, \num{-0.453}
        \unit{PW}; a calota: \qty{-5.28}{PW}.
  \item $\alpha_\text{calota} = \sum\alpha_i\bar Q_iA_i/\sum\bar Q_iA_i = 0.400$, contra \num{0.342}
        do globo.
  \item $-12.90 \times 5.10\times10^{14} = \qty{-6.58}{PW}$: a Terra estaria perdendo calor a essa
        taxa e se resfriando. O desequilíbrio avaliado é $+0.7 \times 5.10\times10^{14} = \qty{+0.36}{PW}$, dezoito vezes menor e de sinal oposto: o desequilíbrio da reanálise é um
        erro do modelo.
  \item Corrigido, $R_i + 12.90$: \num{-15.2}, \num{-38.1}, \num{-63.7}, \num{-92.0},
        \num{-104.0} \unit{W.m^{-2}}; potências \num{-0.478}, \num{-0.972}, \num{-1.197},
        \num{-1.059}, \num{-0.403} \unit{PW}; a calota: \qty{-4.108}{PW}, um déficit médio de
        \qty{45.1}{W.m^{-2}}.
  \item A energia da calota não muda de um ano para o outro, de modo que o que ela perde no topo da
        atmosfera deve entrar nela, e a única maneira de entrar é através do 40.º paralelo:
        $T(\ang{40}) = -2\pi a_\oplus^2\int_{\ang{40}}^{\ang{90}}
        R\cos\varphi\,\mathrm d\varphi$.
  \item $T(\ang{40}\,\text{N}) = +\qty{4.11}{PW}$, para o norte.
  \item $LP = 2.4655\times10^6 \times 1.95/86400 = \qty{55.6}{W.m^{-2}}$, mais que
        $LE = \qty{42.4}{W.m^{-2}}$: cai sobre a calota mais água do que evapora dela; a diferença
        chega como vapor.
  \item $(LE - LP) - 2.07 = 42.4 - 55.6 - 2.07 = \qty{-15.3}{W.m^{-2}}$; vezes a área da calota:
        \qty{-1.40}{PW}. A atmosfera traz \qty{1.40}{PW} de calor latente através de \ang{40} N.
  \item $1.40/4.11 = 0.34$: um terço.
  \item Oceano $0.22 \times 4.11 = \qty{0.90}{PW}$; atmosfera $0.78 \times 4.11 =
        \qty{3.20}{PW}$, dos quais \qty{1.40}{PW} latentes e \qty{1.80}{PW} transportados como calor
        sensível e energia potencial do ar.
  \item Um limite inferior. O total da reanálise a \ang{43} N, \qty{4.02}{PW}, é menor que a
        atmosfera sozinha no estudo, \qty{5.0}{PW}; o total verdadeiro é maior.
  \item Cerca de \qty{4.11}{PW}, \qty{4.11e15}{W}, segundo a radiação da reanálise, um terço dele
        como calor latente; o valor verdadeiro é ainda maior.
        \end{enumerate}
\end{solution}
